\section{Encapsulation}

\begin{frame}{The problem}
Why do we want to make our data attributes protected or private?
\begin{itemize}
    \pause
    \item Sometimes they contain information that shouldn't be known outside of the class
    \pause
    \item \textbf{We want to control how they are changed}
    \begin{itemize}
        \pause
        \item E.g., a student number should always start with a 's' and have 7 numbers.
        \pause
        \item A password should contain certain special characters
    \end{itemize}
    \pause
    \item Sometimes you want a data attribute to be able to be read, but not be able to be set. And the other way around.
\end{itemize}
\end{frame}

\begin{frame}{Solution to this problem}
\textbf{Solution:} we make methods that change our protected and private data attributes (\textbf{setters}) and/or make methods that return the value of our data attribute (\textbf{getters}).
\pause
\begin{itemize}
    \item We '\textbf{encapsulate}' our data attributes with \textit{getters} and/or \textit{setters}.
    \pause
    \item \textbf{This is encapsulation!}
\end{itemize}
\end{frame}

\begin{frame}[plain]
    Let's see how we do this in Python!
\end{frame}

\begin{frame}[fragile]{Making getters and setters}
\begin{minted}[fontsize=\footnotesize]{python}
class Student(Person):
    def __init__(self, first_name, last_name, age, s_number):
        super().__init__(first_name, last_name, age)
        self.__s_number = s_number  # Private data attribute
\end{minted}
\pause
\begin{minted}[fontsize=\footnotesize]{python}
    
    @property  # This is a decorator that indicates its a getter
\end{minted}
\pause
\begin{minted}[fontsize=\footnotesize]{python}
    def s_number(self):  # Note the name
\end{minted}
\pause
\begin{minted}[fontsize=\footnotesize]{python}
        return self.__s_number
\end{minted}
\pause
\begin{minted}[fontsize=\footnotesize]{python}

    @s_number.setter  # Decorator for setter
\end{minted}
\pause
\begin{minted}[fontsize=\footnotesize]{python}
    def s_number(self, new_s_number):  # Note the name
        if new_s_number.startswith("s"):
            self.__s_number = new_s_number
        else:
            raise ValueError
\end{minted}
\end{frame}

\begin{frame}[fragile]{Making getters and setters}
\begin{minted}[fontsize=\footnotesize]{python}
class Student(Person):
    ...
    @property  # This is a decorator that indicates its a getter
    def s_number(self):
        return self.__s_number

    @s_number.setter  # Decorator for setter
    def s_number(self, new_s_number):
        if new_s_number.startswith("s"):
            self.__s_number = new_s_number
        else:
            raise ValueError

student = Student("Daan", "Wichmann", "Secret", "s1234567")
# Because of the 'property' decorator we can use the getter
# as if s_number is just a normal data attribute
print(student.s_number)
>>> "s1234567"
\end{minted}
\end{frame}

\begin{frame}[fragile]{Making getters and setters}
\begin{minted}[fontsize=\footnotesize]{python}
class Student(Person):
    ...
    @property  # This is a decorator that indicates its a getter
    def s_number(self):
        return self.__s_number

    @s_number.setter  # Decorator for setter
    def s_number(self, new_s_number):
        if new_s_number.startswith("s"):
            self.__s_number = new_s_number
        else:
            raise ValueError

student = Student("Daan", "Wichmann", "Secret", "s1234567")
student.s_number = 1234567
print(student.s_number)
\end{minted}
What will happen?
\pause
\begin{minted}[fontsize=\footnotesize]{python}
>>> ValueError
\end{minted}
\end{frame}

\begin{frame}[fragile]{Making getters and setters}
\begin{minted}[fontsize=\footnotesize]{python}
class Student(Person):
    ...
    @property  # This is a decorator that indicates its a getter
    def s_number(self): 
        return self.__s_number

    @s_number.setter  # Decorator for setter
    def s_number(self, new_s_number):
        if new_s_number.startswith("s"):
            self.__s_number = new_s_number
        else:
            raise ValueError

student = Student("Daan", "Wichmann", "Secret", "s1234567")
student.s_number = s1020425
print(student.s_number)
\end{minted}
\pause
\begin{minted}[fontsize=\footnotesize]{python}
>>> "s1020425"
\end{minted}
\end{frame}

\begin{frame}{Encapsulation is important}
\begin{itemize}
    \item Encapsulation is one of the most important concepts of OOP!
    \pause
    \item Make sure to study it from the material, before the exam.
\end{itemize}
\pause
\begin{block}{Encapsulation (definition from the material)}
    Hiding attributes from clients is called encapsulation. As the name implies, the attribute is "enclosed in a capsule". The client is then offered a suitable interface for accessing and processing the data stored in the object.
\end{block}
\pause
\begin{itemize}
    \item Client := your program
    \item The suitable interface that is referred to are the getters and setters
\end{itemize}
\end{frame}